% Part: propositional-logic
% Chapter: syntax-and-semantics
% Section: valuations-sat

\documentclass[../../../include/open-logic-section]{subfiles}

\begin{document}

\olfileid[hi]{pl}{syn}{val}

\olsection{\usetoken{S}{valuation} और संतुष्टि}

\begin{defn}[!!^{valuation}]
$\{\True, \False\}$ दो सत्य-मूल्यों, ``सत्य'' और ``असत्य'', का समुच्चय हो।
$\Lang{L_0}$ के लिए \emph{!!{valuation}} वह फलन~$\pAssign{v}$ है जो भाषा के
!!{propositional variable}s में से प्रत्येक को $\True$ या $\False$ देता है;
अर्थात् $\pAssign{v} \colon
\PVar \to \{\True, \False \}$।
\end{defn}

\begin{defn}
  दिए गए !!a{valuation}~$\pAssign{v}$ के लिए मूल्य-फलन
  $\pValue{v} \colon \Frm[L_0] \to \{\True, \False \}$ को निम्न प्रकार
  आगमनात्मक रूप से परिभाषित करें:
  \begin{align*}
    \iftag{prvFalse}{\pValue{v}(\lfalse) & = \False; \\}{}
    \iftag{prvTrue}{\pValue{v}(\ltrue) & = \True; \\}{}
    \pValue{v}(\Obj p_n) & = \pAssign{v}(\Obj p_n); \\
    \iftag{prvNot}{\pValue{v}(\lnot !A) & = \begin{cases}
        \True & \text{यदि } \pValue{v}(!A) = \False;\\
        \False & \text{अन्यथा।}
      \end{cases} \\ }{}
    \iftag{prvAnd}{\pValue{v}(!A \land !B) & = \begin{cases}
        \True &
        \text{यदि $\pValue{v}(!A) = \True$ और $\pValue{v}(!B) = \True$;}\\
        \False &
        \text{यदि $\pValue{v}(!A) = \False$ या $\pValue{v}(!B) = \False$}.
      \end{cases}\\}{}
    \iftag{prvOr}{\pValue{v}(!A \lor !B) & = \begin{cases}
        \True &
        \text{यदि $\pValue{v}(!A) = \True$ या $\pValue{v}(!B) = \True$;}\\
        \False &
        \text{यदि $\pValue{v}(!A) = \False$ और $\pValue{v}(!B) = \False$}.
      \end{cases}\\}{}
    \iftag{prvIf}{\pValue{v}(!A \lif !B) & = \begin{cases}
        \True &
        \text{यदि $\pValue{v}(!A) = \False$ या $\pValue{v}(!B) = \True$;}\\
        \False &
        \text{यदि $\pValue{v}(!A) = \True$ और $\pValue{v}(!B) = \False$}.
      \end{cases}\\}{}
    \iftag{prvIff}{\pValue{v}(!A \liff !B) & = \begin{cases}
        \True &
        \text{यदि $\pValue{v}(!A) = \pValue{v}(!B)$;}\\
        \False &
        \text{यदि $\pValue{v}(!A) \neq \pValue{v}(!B)$}.
      \end{cases}}{}
\end{align*}
\end{defn}

\begin{explain}
ये खंड निम्नलिखित सत्यता-सारणियों के अनुरूप हैं:
\begin{center}
\iftag{prvNot}{
  \begin{tabular}{|c||c|} \hline
    $!A$ & $ \lnot !A$ \\
    \hline \hline
    $\True$ & $\False$ \\
    $\False$ & $\True$ \\
    \hline
  \end{tabular}
}{}
\iftag{prvAnd}{
  \begin{tabular}{|cc||c|} \hline
    $!A$ & $!B$ & $!A \land !B$ \\
    \hline \hline
    $\True$ & $\True$ & $\True$ \\
    $\True$ & $\False$ & $\False$ \\
    $\False$ & $\True$ & $\False$ \\
    $\False$ & $\False$ & $\False$ \\
    \hline
  \end{tabular}
}{}
\iftag{prvOr}{
  \begin{tabular}{|cc||c|} \hline
    $!A$ & $!B$ & $!A \lor !B$ \\
    \hline \hline
    $\True$ & $\True$ & $\True$ \\
    $\True$ & $\False$ & $\True$ \\
    $\False$ & $\True$ & $\True$ \\
    $\False$ & $\False$ & $\False$ \\
    \hline
  \end{tabular}
}{}%
\iftag{notprvNot,notprvAnd,notprvOr,notprvIf}{}{\\[1em]}
\iftag{prvIf}{
  \begin{tabular}{|cc||c|} \hline
    $!A$ & $!B$ & $!A \lif !B$ \\
    \hline \hline
    $\True$ & $\True$ & $\True$ \\
    $\True$ & $\False$ & $\False$ \\
    $\False$ & $\True$ & $\True$ \\
    $\False$ & $\False$ & $\True$ \\
    \hline
  \end{tabular}
}{}
\iftag{prvIff}{
  \begin{tabular}{|cc||c|} \hline
    $!A$ & $!B$ & $!A \liff !B$ \\
    \hline \hline
    $\True$ & $\True$ & $\True$ \\
    $\True$ & $\False$ & $\False$ \\
    $\False$ & $\True$ & $\False$ \\
    $\False$ & $\False$ & $\True$ \\
    \hline
  \end{tabular}
}{}
\end{center}
\end{explain}

\begin{prob}
$\Lang{L_0}$ में एक त्रिस्थानी संयोजक $\diamondsuit$ जोड़ने पर विचार करें,
जिसका मूल्यांकन इस प्रकार दिया गया है:
\begin{gather*}
  \pValue{v}(\diamondsuit( !A , !B , !C ) ) = \begin{cases}
    \pValue{v}( !B ) &
    \text{यदि $\pValue{v}(!A) = \True$;}\\
    \pValue{v}( !C ) &
    \text{यदि $\pValue{v}(!A) = \False $}.
  \end{cases}
\end{gather*}
इस संयोजक की सत्यता-सारणी लिखें।
\end{prob}

\begin{thm}[स्थानीय निर्धारण]
\ollabel{thm:LocalDetermination} मानें कि $\pAssign{v_1}$ और
$\pAssign{v_2}$ !!{valuation} हैं और $!A$ में आने वाले
!!{propositional
variable}s पर सहमत हैं; अर्थात् $\pAssign{v_1}(\Obj p_n) =
\pAssign{v_2}(\Obj p_n)$, जब भी $\Obj p_n$ किसी
!!{formula}~$!A$ में आए। तब $\pValue{v_1}$ और $\pValue{v_2}$ भी~$!A$ पर
सहमत हैं; अर्थात् $\pValue{v_1}(!A) = \pValue{v_2}(!A)$।
\end{thm}

\begin{proof}
$!A$ पर आगमन द्वारा।
\end{proof}

\begin{defn}[संतुष्टि]
\ollabel{defn:satisfaction} हम \emph{!!a{formula}~$!A$ की
  !!a{valuation}~$\pAssign{v}$ द्वारा संतुष्टि} की धारणा $\pSat{v}{!A}$ को
  निम्न प्रकार आगमनात्मक रूप से परिभाषित कर सकते हैं।
  (हम $\pSat/{v}{!A}$ को ``$\pSat{v}{!A}$ नहीं'' के अर्थ में लिखते हैं।)
\begin{enumerate}
\tagitem{prvFalse}{%
  \indcase{!A}{\lfalse}{$\pSat/{v}{\indfrm}$।}}{}

\tagitem{prvTrue}{%
  \indcase{!A}{\ltrue}{$\pSat{v}{\indfrm}$।}}{}

\item \indcase{!A}{\Obj p_i}{$\pSat{v}{\indfrm}$
  तभी और केवल तभी, जब $\pAssign{v}(\Obj p_i) = \True$।}

\tagitem{prvNot}{%
  \indcase{!A}{\lnot !B}{$\pSat{v}{\indfrm}$ तभी और केवल तभी, जब
    $\pSat/{v}{!B}$।}}{}

\tagitem{prvAnd}{%
  \indcase{!A}{(!B \land !C)}{$\pSat{v}{\indfrm}$ तभी और केवल तभी, जब
    $\pSat{v}{!B}$ और $\pSat{v}{!C}$।}}{}

\tagitem{prvOr}{%
  \indcase{!A}{(!B \lor !C)}{$\pSat{v}{\indfrm}$ तभी और केवल तभी, जब
    $\pSat{v}{!B}$ या $\pSat{v}{!C}$ (या दोनों)।}}{}

\tagitem{prvIf}{%
  \indcase{!A}{(!B \lif !C)}{$\pSat{v}{\indfrm}$ तभी और केवल तभी, जब
    $\pSat/{v}{!B}$ या $\pSat{v}{!C}$ (या दोनों)।}}{}

\tagitem{prvIff}{%
  \indcase{!A}{(!B \liff !C)}{$\pSat{v}{\indfrm}$ तभी और केवल तभी, जब या तो
    $\pSat{v}{!B}$ और $\pSat{v}{!C}$ दोनों, अथवा न $\pSat{v}{!B}$
    न $\pSat{v}{!C}$।}}{}
\end{enumerate}
यदि $\Gamma$ !!{formula}s का कोई समुच्चय है, तो $\pSat{v}{\Gamma}$ तभी और
केवल तभी, जब $\pSat{v}{!A}$ प्रत्येक~$!A \in \Gamma$ के लिए सत्य हो।
\end{defn}

\begin{prop}\ollabel{prop:sat-value}
  $\pSat{v}{!A}$ तभी और केवल तभी, जब $\pValue{v}(!A) = \True$।
\end{prop}

\begin{proof}
  $!A$ पर आगमन द्वारा।
\end{proof}

\begin{prob}
  \olref[pl][syn][val]{prop:sat-value} सिद्ध करें।
\end{prob}

\end{document}
